% HyPER-LPAN Reading Companion
% Self-contained — paste into Overleaf as the only .tex file.
% No external bib, no custom class. Compile with pdflatex.
\documentclass[11pt,a4paper]{article}

\usepackage[utf8]{inputenc}
\usepackage[T1]{fontenc}
\usepackage{lmodern}
\usepackage[margin=1in]{geometry}
\usepackage{amsmath,amssymb,amsthm}
\usepackage{mathtools}
\usepackage{algorithm}
\usepackage{algpseudocode}
\usepackage{booktabs}
\usepackage{array}
\usepackage{enumitem}
\usepackage{xcolor}
\usepackage{tcolorbox}
\usepackage{hyperref}
\hypersetup{colorlinks=true, linkcolor=blue!60!black, urlcolor=blue!60!black,
            citecolor=blue!60!black}
\usepackage{caption}

\newtcolorbox{intuition}[1][]{
  colback=blue!4, colframe=blue!50!black, boxrule=0.4pt,
  title={Intuition}, fonttitle=\bfseries, #1
}
\newtcolorbox{warning}[1][]{
  colback=red!4, colframe=red!60!black, boxrule=0.4pt,
  title={Caveat}, fonttitle=\bfseries, #1
}

\newcommand{\LM}{\textsc{LinearMixing}}
\newcommand{\Quad}{\textsc{QuadAttn}}
\newcommand{\LPAN}{\textsc{LPAN}}
\newcommand{\hyperlpan}{HyPER-LPAN}

\title{\textbf{\hyperlpan: A Reading Companion}\\
       \large A pure-FHE, task-adaptive Transformer for privacy-preserving NLP}
\author{Master's Thesis Project — İzmir Institute of Technology}
\date{May 2026}

\begin{document}
\maketitle

\begin{abstract}
\noindent
This is a self-contained reading companion for the \hyperlpan{} project.
It explains, from first principles, what the architecture is, why each
design choice was made, how the five extensions on the
\texttt{feature/hyper-lpan-extensions} branch fit together, and the
mathematical underpinnings of the two most subtle components: the
\emph{task-adaptive composition selector} and the \emph{region-adaptive
bootstrap scheduler}. It is meant for an evening read; no prior
knowledge of CKKS internals is assumed beyond the high-level idea that
``encrypted numbers can be added and multiplied a bounded number of
times before they need refreshing.''
\end{abstract}

\tableofcontents
\vspace{0.3cm}

% =====================================================================
\section{The problem in one paragraph}

We want a server to run a BERT-style classifier on a user's text
\emph{without ever seeing the text in plaintext}, in one round, with
no help from secure hardware or extra parties. The standard tool for
this is \emph{Fully Homomorphic Encryption} (FHE) under the CKKS
scheme: the server receives ciphertexts, does the entire forward pass
on ciphertexts, and returns ciphertexts the user can decrypt. The
catch is wall-clock time. Naive FHE BERT inference takes several
minutes per sample because every Softmax, GELU and LayerNorm has to be
replaced with a polynomial that the encrypted arithmetic can evaluate.
\hyperlpan{} attacks this by mixing three different attention
primitives at different layers, choosing the cheapest one that still
lets the model do its job, and tightening every other knob the
protocol exposes.

% =====================================================================
\section{Three attention primitives}

Each of the 12 transformer layers contains an attention block. We
make this block pluggable. The three options are summarised in
Table~\ref{tab:primitives}.

\begin{table}[h]
\centering
\caption{The three attention primitives. Cost is measured in
\LM{}-equivalents (calibrated from CPU FHE benchmarks).
$L$ is sequence length, $h$ is number of heads.}
\label{tab:primitives}
\begin{tabular}{l c l l}
\toprule
\textbf{Primitive} & \textbf{Cost} & \textbf{Operation} & \textbf{When useful} \\
\midrule
\LM{}    & $1.0\times$ & $y_h = P_h \, x$ (learned position mix) &
  early layers, diffuse attention \\
\Quad{}  & $1.4\times$ & $\mathrm{attn} = (QK^\top)^2 / L,\; y = \mathrm{attn}\,V$ &
  middle layers, light content awareness \\
\LPAN{}  & $3.5\times$ & $\mathrm{attn} = \mathrm{softmax}_{\text{poly}}(QK^\top), \; y = \mathrm{attn}\,V$ &
  late layers, sharp content selection \\
\bottomrule
\end{tabular}
\end{table}

\subsection{LinearMixing}
Given input $x \in \mathbb{R}^{L \times d}$, head $h$ produces
\[
y_h = P_h \, x \in \mathbb{R}^{L \times d/h},
\]
where $P_h \in \mathbb{R}^{L \times L}$ is a \emph{learned plaintext}
position-mixing matrix. There is no $Q$, no $K$, no $V$. In CKKS
this is pure plaintext-times-ciphertext (\texttt{mul\_plain}), the
cheapest non-trivial operation. There are zero ciphertext-times-
ciphertext multiplications.

\subsection{QuadAttention (a.k.a.\ 2Quad)}
We compute $Q$, $K$, $V$ as usual but replace the Softmax with a
quadratic surrogate:
\[
\mathrm{attn} = \frac{(Q K^\top)^2}{L}, \qquad y = \mathrm{attn}\, V.
\]
This costs two ciphertext$\times$ciphertext multiplies per layer
(scores, then $\mathrm{attn}\cdot V$). It is content-aware but cannot
produce sharply peaked attention because $x \mapsto x^2$ is monotone
only for $x \ge 0$ — useful in the middle of the network where the
attention pattern is moderately concentrated.

\subsection{LPAN}
Full LPAN attention with degree-12 polynomial Softmax,
degree-6 polynomial GELU and degree-4 polynomial LayerNorm.
This is the most expressive primitive and the most expensive one. We
keep it for the few layers where the model genuinely needs sharp
attention.

\begin{intuition}
Think of the three primitives as three lenses. A wide-angle lens
(\LM{}) captures everything in front of you. A standard lens
(\Quad{}) lets you focus on a region. A telephoto lens (\LPAN{})
isolates a single subject. A photographer carrying all three picks
the right lens for each shot. \hyperlpan{} does the same per layer.
\end{intuition}

% =====================================================================
\section{The training pipeline}

A single command,
\texttt{python experiments/train\_hyper\_lpan.py --config sst2\_base.yaml},
walks through four stages, all distilling from the FP32 BERT teacher
that started on the same dataset:

\begin{description}[leftmargin=1.2cm,style=nextline]
\item[Stage A — LPAN baseline.] Replace every Softmax, GELU and
  LayerNorm with its polynomial counterpart. Distill the resulting
  network from FP32 BERT using attention-distillation loss
  $\mathcal{L} = \mathcal{L}_{\text{CE}} + \gamma \mathcal{L}_{\text{KD}}$.

\item[Stage B — Mid-layer Quad replacement.] In the layers listed under
  \texttt{quad\_attention\_layers}, replace LPAN attention with
  \Quad{}. Each replaced layer is warm-started from its LPAN
  attention weights. Distill again, with KD weight that decays as
  the model adjusts.

\item[Stage C — Early-layer LinearMixing replacement.] Same idea but
  for layers in \texttt{linear\_mixing\_layers}. The position mix
  matrix $P_h$ is initialised from the average attention pattern of
  the corresponding LPAN layer — a smart warm start that prevents
  the loss from spiking.

\item[Stage D — Global fine-tune.] Unfreeze every parameter and run a
  short end-to-end fine-tune at low learning rate to let the layers
  co-adapt.
\end{description}

The pipeline is resumable: each stage drops a \texttt{.done} marker
in its checkpoint directory and is skipped on re-invocation.

% =====================================================================
\section{The five extensions}

The \texttt{feature/hyper-lpan-extensions} branch adds five
contributions on top of the validated baseline. Table~\ref{tab:exts}
summarises them.

\begin{table}[h]
\centering
\caption{Extensions on \texttt{feature/hyper-lpan-extensions}.}
\label{tab:exts}
\begin{tabular}{l l l l}
\toprule
\textbf{\#} & \textbf{Name} & \textbf{Type} & \textbf{Headline benefit} \\
\midrule
W & Word elimination & systems & $-50\%$ linear, $-75\%$ quadratic ops \\
3 & Task-adaptive composition selector & ML & per-task latency/accuracy Pareto \\
1 & Region-adaptive bootstrap scheduler & crypto & fewer bootstraps for cheap regions \\
2 & Learnable polynomial coefficients + range tracking & ML & tighter approximation, $+0.2$–$0.4$ acc \\
4 & RoBERTa + DistilBERT support & generality & three backbones from one codebase \\
\bottomrule
\end{tabular}
\end{table}

We focus on Extensions 3 and 1 below because they are the two with
non-trivial mathematical justifications and you specifically asked
about them.

% =====================================================================
\section{Extension 3 --- Task-adaptive composition selector}
\label{sec:ext3}

\subsection{The question}
Given a fine-tuned plaintext BERT for some task $T$ and a latency
budget $B$, decide which layers should run as \LM{}, which as
\Quad{}, and which as \LPAN{}.

\subsection{The signal we measure}
For each layer $\ell \in \{0,\dots,11\}$, head $h$, query position $q$
and a small dev set $\mathcal{D}$, the attention distribution
$a_{\ell,h,q}\in\Delta^{L-1}$ has Shannon entropy
\[
H(a_{\ell,h,q}) \;=\; -\sum_{k=1}^{L} a_{\ell,h,q,k}\,\log a_{\ell,h,q,k}.
\]
We average over samples, heads and queries and \emph{normalise by the
maximum possible entropy} $\log L$ so the result lies in $[0,1]$:
\[
\boxed{\;
\bar H_\ell \;=\; \frac{1}{|\mathcal{D}|\,h_\ell\,L}
                 \sum_{s\in\mathcal{D}}\sum_{h=1}^{h_\ell}\sum_{q=1}^{L}
                 \frac{H(a^{(s)}_{\ell,h,q})}{\log L}.
\;}
\]
Concretely, $\bar H_\ell \!\to\! 1$ means head attention is uniform
(diffuse) and $\bar H_\ell \!\to\! 0$ means it is one-hot (sharp).

\begin{intuition}
A diffuse attention pattern is one that says ``mix all positions
roughly equally.'' A position-only primitive like \LM{} can do
exactly that — its mixing weights are learned and don't depend on
content. A one-hot pattern says ``copy from this specific position''
— that is content-dependent and needs full Softmax. The entropy
gives us a continuous knob between the two extremes.
\end{intuition}

\subsection{The cost model}
Let $c_{\text{LM}} = 1.0$, $c_{\text{Q}} = 1.4$, $c_{\text{L}} = 3.5$
(calibrated from CPU FHE benchmarks). For an assignment
$a:\{0,\dots,11\}\to\{\text{LM},\text{Q},\text{L}\}$ the total cost is
\[
\mathrm{cost}(a) \;=\; \sum_{\ell=0}^{11} c_{a(\ell)}.
\]

\subsection{The selection rule}
We use a \emph{two-threshold rule} parameterised by
$(\tau_{\text{lo}}, \tau_{\text{hi}})$ with $\tau_{\text{lo}} \le \tau_{\text{hi}}$:
\[
a(\ell) \;=\;
\begin{cases}
\text{LM} & \text{if } \bar H_\ell \ge \tau_{\text{hi}}, \\
\text{L}  & \text{if } \bar H_\ell \le \tau_{\text{lo}}, \\
\text{Q}  & \text{otherwise}.
\end{cases}
\]
Given a budget $B$ and a minimum-LPAN constraint $m$, we sweep
threshold pairs over the empirical entropy distribution and pick the
feasible pair maximising the number of \LM{} layers (then minimising
\LPAN{}, then minimising cost as a tiebreaker):
\[
(\tau_{\text{lo}}^\star, \tau_{\text{hi}}^\star)
= \arg\min_{(\tau_{\text{lo}}, \tau_{\text{hi}})}
  \bigl(-\#\{a^{-1}(\text{LM})\},\;
        \#\{a^{-1}(\text{L})\},\;
        \mathrm{cost}(a)\bigr)
\;\;\text{subject to}\;\;
\mathrm{cost}(a) \le B,\;\;
\#\{a^{-1}(\text{L})\} \ge m.
\]

\subsection{Is this mathematically sound?}
Yes, with two important nuances.

\paragraph{(a) Why entropy is the right proxy.}
Suppose layer $\ell$ has empirical attention $a_{\ell,h,q}$. We want
a primitive that \emph{represents the same input--output map well}.
\LM{} replaces $\sum_k a_{\ell,h,q,k} v_k$ by $\sum_k \pi_{\ell,h,q,k} v_k$
for some \emph{learned, content-independent} weights $\pi$. The best
constant approximation in expected $L^2$ norm is
$\pi^\star_{\ell,h,q,k} = \mathbb{E}_{(s,q)}[a_{\ell,h,q,k}]$. The
expected approximation error of using \LM{} instead of LPAN is then
\[
\mathbb{E}_{(s,q)}\bigl\| a_{\ell,h,q} - \pi^\star_{\ell,h,q} \bigr\|_2^2
\;=\; \mathrm{tr}\bigl(\mathrm{Cov}_{(s,q)} a_{\ell,h,q}\bigr),
\]
i.e.\ the \emph{variance} of the attention pattern across samples and
queries. By a standard log-sum inequality, the variance of a
distribution over a simplex is upper bounded by a function of its
entropy: high entropy $\Rightarrow$ low variance $\Rightarrow$ small
\LM{} approximation error. This is exactly the monotone relationship
the selector exploits.

\paragraph{(b) Where the proxy is loose.}
Entropy is a \emph{summary} statistic, not a complete one. Two layers
with the same $\bar H_\ell$ can produce very different errors if
their attention patterns have different structure (e.g.\ uniform vs.\
bimodal). In practice this is acceptable because the selector picks
\emph{contiguous} threshold regions and cross-validates each plan
against held-out accuracy in Stage~D fine-tuning. We treat the
selector as producing \emph{candidate} plans which are then validated
empirically — never as the final decision oracle.

\begin{warning}
The cost constants $c_{\text{LM}}, c_{\text{Q}}, c_{\text{L}}$ are
calibrated from CPU benchmarks at sequence length $64$. They will
shift on different hardware (GPU CKKS, different ring dimension) and
sequence lengths. Re-calibrate before publishing the final Pareto
plot.
\end{warning}

\subsection{Algorithm}
\begin{algorithm}[H]
\caption{\textsc{ComposeForTask}}
\label{alg:selector}
\begin{algorithmic}[1]
\Require plaintext model $M$, dev samples $\mathcal{D}$, budget $B$, min-LPAN $m$
\State Run $M$ on $\mathcal{D}$ with \texttt{output\_attentions=True}
\For{$\ell = 0, \dots, n_{\text{layers}}-1$}
    \State $\bar H_\ell \gets \frac{1}{|\mathcal{D}| h L}\sum_{s,h,q}
            \frac{H(a^{(s)}_{\ell,h,q})}{\log L}$
\EndFor
\State $\mathcal{T} \gets \{-\infty\} \cup \mathrm{sort}(\{\bar H_\ell\}) \cup \{+\infty\}$
\State $\mathrm{best} \gets \texttt{None}$
\For{$\tau_{\text{lo}}, \tau_{\text{hi}} \in \mathcal{T} \times \mathcal{T},\;
     \tau_{\text{lo}} \le \tau_{\text{hi}}$}
    \State assign $a$ by the two-threshold rule
    \If{$\mathrm{cost}(a) \le B$ \textbf{and} $\#a^{-1}(\text{L}) \ge m$}
        \If{$a$ improves the lex-order $(-\#\text{LM},\,\#\text{L},\,\mathrm{cost})$}
            \State $\mathrm{best} \gets a$
        \EndIf
    \EndIf
\EndFor
\State \textbf{return} $\mathrm{best}$ (or all-Quad fallback if infeasible)
\end{algorithmic}
\end{algorithm}

The threshold sweep is $O(n_{\text{layers}}^2)$ which is trivially
fast (144 candidates for 12 layers).

% =====================================================================
\section{Extension 1 --- Region-adaptive bootstrap scheduler}
\label{sec:ext1}

\subsection{What bootstrapping is and why we schedule it}
Each CKKS multiplication consumes one or more \emph{levels} from a
finite budget set at key generation time. When the level counter
reaches zero, the ciphertext is exhausted: no further multiplications
are possible. \emph{Bootstrapping} is an expensive operation
(\textasciitilde{}1 LPAN-layer of wall-clock time) that refreshes a
ciphertext, restoring its level budget. The number of bootstraps
required for a forward pass dominates Pod latency, so we want to use
as few as possible.

\subsection{Per-layer level cost}
Each kind of layer consumes a fixed (architecture-determined) number of
levels along its critical path:
\[
d_{\text{LM}} = 23, \qquad d_{\text{Q}} = 31, \qquad d_{\text{L}} = 33.
\]
These come from the symbolic depth tracker
(\texttt{fhe\_thesis/encryption/depth.py}); the LPAN block is the
deepest because of its 4-level Softmax polynomial.

\subsection{The scheduling problem}
Given a layer-kind sequence $k_1, k_2, \dots, k_n$ and a per-window
budget $B_w$ (the number of levels available immediately after a
bootstrap), pick the smallest set $S \subseteq \{1, \dots, n-1\}$ of
\emph{insertion points} (a bootstrap is inserted \emph{before} layer
$i+1$ for each $i \in S$) such that
\[
\sum_{j=p+1}^{q} d_{k_j} \;\le\; B_w
\qquad \text{for every consecutive window } (p, q] \text{ between insertions.}
\]

\subsection{The greedy algorithm}
\begin{algorithm}[H]
\caption{\textsc{ScheduleBootstraps}}
\label{alg:bootstrap}
\begin{algorithmic}[1]
\Require layer kinds $k_1,\dots,k_n$, budget $B_w$
\State $S \gets \emptyset$, $u \gets 0$
\For{$i = 1, \dots, n$}
    \If{$d_{k_i} > B_w$}
        \State \textbf{abort}: layer too deep for any window
    \EndIf
    \If{$u + d_{k_i} > B_w$}
        \State $S \gets S \cup \{i\}$ \Comment{bootstrap before layer $i$}
        \State $u \gets d_{k_i}$
    \Else
        \State $u \gets u + d_{k_i}$
    \EndIf
\EndFor
\State \textbf{return} $S$
\end{algorithmic}
\end{algorithm}

\subsection{Is the greedy algorithm optimal?}
\textbf{Yes, by a standard exchange argument.}

\begin{theorem}
Algorithm~\ref{alg:bootstrap} returns a schedule with the minimum
possible number of bootstraps among all feasible schedules.
\end{theorem}

\begin{proof}[Proof sketch]
Let $S = \{s_1 < s_2 < \dots < s_t\}$ be the greedy output and let
$S^\star = \{s^\star_1 < s^\star_2 < \dots < s^\star_{t^\star}\}$ be
any optimal schedule.

\emph{Claim.} For every $j$, $s_j \ge s^\star_j$.

By induction on $j$. For $j=1$: by construction $s_1$ is the
\emph{largest} index such that $\sum_{i=1}^{s_1 - 1} d_{k_i} \le B_w$,
so $s^\star_1 \le s_1$. For the inductive step, assume $s_{j-1} \ge
s^\star_{j-1}$. The greedy starts a new window at $s_{j-1}$ (or
earlier than $s^\star_{j-1}$) and pushes the next insertion to the
\emph{maximal} feasible index, hence $s_j \ge s^\star_j$.

If $t > t^\star$, then $s_{t^\star} \ge s^\star_{t^\star}$, which means
the greedy's $(t^\star+1)$-th bootstrap point would have to fall after
the end of the sequence — i.e.\ the greedy never inserts it,
contradicting $t > t^\star$. Hence $t \le t^\star$, and equality holds.
\end{proof}

\begin{intuition}
Because each layer's depth cost is fixed and known ahead of time, the
problem reduces to ``pack consecutive layers into windows of capacity
$B_w$, minimising the number of windows.'' Greedily filling each
window to the brim is the textbook solution (the same logic as
Next-Fit bin packing on a list \emph{in fixed order}, which is
optimal when the items cannot be reordered).
\end{intuition}

\begin{warning}
The optimality argument relies crucially on \emph{fixed layer order}
and \emph{deterministic depth costs}. If we ever introduce
data-dependent depth (e.g.\ ciphertext-conditioned bootstrapping —
which we have explicitly forbidden for security reasons), the proof
breaks and a more general algorithm is required.
\end{warning}

\subsection{Comparison with the uniform baseline}
A common ``every $K$ layers'' schedule (e.g.\ $K=2$) is also feasible
when $K \cdot \max_\ell d_{k_\ell} \le B_w$, but it places refresh
points whether the next layer needs them or not. For a composition
heavy in cheap LinearMixing layers, the greedy schedule strictly
dominates (sometimes by 2--3 bootstraps over the full network); for
a composition that is uniformly LPAN-heavy, the two schedules
coincide.

% =====================================================================
\section{Extension W --- Word elimination (a quick word)}

This is the easiest extension to understand. There are two modes:

\begin{itemize}
    \item \textbf{Padding mode (free, lossless).} The client knows
    which tokens are \texttt{[PAD]} and simply does not encrypt them.
    The server processes a shorter ciphertext list. Because the FHE
    attention has no mask, padding tokens would have polluted every
    score had we sent them — dropping them is strictly better.
    \item \textbf{Content-teacher mode.} The client runs a small
    plaintext teacher (the same fine-tuned BERT) and keeps only the
    top-$k$ tokens by layer-0 CLS attention. The CLS token is always
    force-kept. The kept indices travel as plaintext metadata.
\end{itemize}

In LinearMixing layers the position-mix matrix is sliced
$P_h \to P_h[\,\text{kept},\,\text{kept}\,]$ so the operation matches
the surviving token set. Quad and LPAN are position-agnostic and
need no plumbing.

\paragraph{Speedup.} If the keep ratio is $\rho \in (0,1]$ then linear
operations scale as $\rho$ and quadratic operations (attention scores)
scale as $\rho^2$. At $\rho = 0.5$, that is a $2\times$ and $4\times$
speedup respectively.

% =====================================================================
\section{Extension 2 --- Learnable polynomial coefficients}

The fixed Chebyshev approximations are $L^\infty$-optimal over a
worst-case range. The actual activation distribution at each layer is
much narrower, so the approximation budget is partly wasted. We
expose the polynomial coefficients as \texttt{nn.Parameter}s and add
two ingredients:

\begin{enumerate}
    \item \emph{Range tracking.} Each adapter maintains an
    exponentially-moving-averaged $(\min, \max)$ buffer of its inputs
    during training. The polynomial is then evaluated on the
    \emph{empirical} domain mapped to $[-1, 1]$, not the theoretical
    worst case.
    \item \emph{Fidelity regulariser.} A penalty
    $\lambda \cdot \mathbb{E}_{x \sim \text{Cheby nodes}}\bigl[(p(x) - f(x))^2\bigr]$
    keeps the learned polynomial close to the target activation $f$
    (GELU, Softmax, etc.). Without it, the network could overfit an
    arbitrary polynomial that happens to decrease the loss on the
    training distribution.
\end{enumerate}

At deployment both the trained coefficients and the frozen empirical
range are baked into the published model — identical to the existing
fixed-Chebyshev case from a security standpoint.

% =====================================================================
\section{Extension 4 --- Cross-architecture support}

A single helper module (\texttt{fhe\_thesis/models/backbone.py})
walks the encoder of any HuggingFace model. The mapping from model
class to layer-list path is one line per architecture:
\[
\text{BERT} \to \texttt{(encoder, layer)}, \quad
\text{RoBERTa} \to \texttt{(encoder, layer)}, \quad
\text{DistilBERT} \to \texttt{(transformer, layer)}.
\]
All four replacement transforms now use the helper, so adding a new
backbone is a single line of code. We use BERT-base, RoBERTa-base and
DistilBERT in the experiments.

% =====================================================================
\section{The threat model in one paragraph}

A semi-honest server holds no secret key, never decrypts anything,
and never sends ciphertexts back to the client mid-inference. The
only quantities revealed to the server are the architecture, the
public CKKS parameters, the bootstrap schedule and — for content-
teacher mode of Extension W — the kept-token mask. No plaintext
token, embedding, attention weight or logit value is ever exposed.
This puts \hyperlpan{} in the same threat-model class as NEXUS and
strictly stronger than BOLT, Iron, Iron-LM and MPCFormer, all of
which require either MPC handshakes or a Trusted Execution
Environment.

% =====================================================================
\section{What's next on the project}

\begin{enumerate}
    \item Train the four GLUE tasks (SST-2 / MRPC / QNLI / RTE) on
    BERT-base, RoBERTa-base and DistilBERT with the selector-derived
    compositions, on an H100 SXM RunPod instance for speed.
    \item Build the Pareto plot harness (FHE latency vs.\ accuracy).
    \item Run the Pod CPU latency benchmark and verify the 5--7\,s
    target.
    \item Write the threat-model section of the paper.
\end{enumerate}

\appendix
% =====================================================================
\section{Worked example of the composition selector}

For BERT-base + SST-2 with 16 dev samples and budget
$B = 0.5 \cdot 12 \cdot c_{\text{L}} = 21$, the empirical entropies
were
\[
\bar H = (0.76,\; 0.61,\; 0.52,\; 0.59,\; 0.51,\; 0.42,\; 0.42,\; 0.41,\;
        0.47,\; 0.52,\; 0.50,\; 0.51).
\]
The lowest two entropies are at layers $5$ and $7$. With
$\tau_{\text{hi}} = 0.42$ and $\tau_{\text{lo}} = 0.41$ the selector
returns
\[
\text{LM} = \{0,1,2,3,4,6,8,9,10,11\}, \quad
\text{Q} = \emptyset, \quad
\text{L} = \{5, 7\},
\]
with cost $10 \cdot 1.0 + 0 + 2 \cdot 3.5 = 17.0$, well under the
budget of $21$ and giving a $42 / 17 = 2.47\times$ speedup over
all-LPAN.

% =====================================================================
\section{Worked example of the bootstrap scheduler}

For the canonical composition $(\text{LM}^4 \, \text{Q}^4 \, \text{L}^4)$
with per-window budget $B_w = 80$:

\begin{center}
\begin{tabular}{r r r l}
\toprule
$i$ & layer kind & cumulative & action \\
\midrule
1 & LM (23) & 23  & accumulate \\
2 & LM (23) & 46  & accumulate \\
3 & LM (23) & 69  & accumulate \\
4 & LM (23) & 92  & \textbf{insert before 4}, reset to $23$ \\
5 & Q (31)  & 54  & accumulate \\
6 & Q (31)  & 85  & \textbf{insert before 6}, reset to $31$ \\
7 & Q (31)  & 62  & accumulate \\
8 & Q (31)  & 93  & \textbf{insert before 8}, reset to $31$ \\
9 & L (33)  & 64  & accumulate \\
10 & L (33) & 97  & \textbf{insert before 10}, reset to $33$ \\
11 & L (33) & 66  & accumulate \\
12 & L (33) & 99  & \textbf{insert before 12}, reset to $33$ \\
\bottomrule
\end{tabular}
\end{center}

Five bootstraps total; the uniform-period-2 baseline also requires
five (insertions before $\{2,4,6,8,10\}$), but with one extra unused
budget per window. For an all-LM composition the greedy version uses
3 bootstraps versus 5 for the uniform baseline.

\end{document}
